\documentclass[11pt]{article}
\usepackage[a4paper,margin=24mm]{geometry}
\usepackage[T1]{fontenc}
\usepackage{lmodern,amsmath,amssymb,amsthm,mathtools,microtype,enumitem,booktabs}
\usepackage[colorlinks=true,linkcolor=blue,urlcolor=blue]{hyperref}
\hypersetup{pdftitle={The operator picture of the proof that zeta(3) is irrational},pdfauthor={Mathathon peer-review digest}}
\newcommand{\Z}{\mathbb Z}
\newcommand{\ip}[2]{\langle #1,#2\rangle}
\newcommand{\q}[1]{\medskip\noindent\textbf{#1}}
\setlist{itemsep=1pt,topsep=3pt}
\title{The operator picture of the proof that $\zeta(3)$ is irrational\\[2pt]
\large what the language of Hamiltonians and commutators says, and what it buys}
\author{Mathathon peer-review digest (Claude)}
\date{14 November 2026}
\begin{document}
\maketitle

\noindent This note collects what four AI agents and I found when we asked how the ``Hamiltonian''
language of the Codex kernel-energy proof can be used. It is a map, not a new proof. The short
answer: there is one operator $H$, and
\begin{center}
\emph{$H$ is the chain of resistors $1/n^3$ whose total resistance is $\zeta(3)$, written as a
differential operator; it is also the Laplacian of the hyperbolic plane.}
\end{center}
Everything in the sum proof (the companion note \texttt{max-principle-proof.pdf}, cited as [MP]) is a
statement about $H$ in the basis of Legendre polynomials. Each statement below carries a label:
\textbf{[V]} verified by me symbolically or in exact arithmetic, \textbf{[A]} reported by an agent
and not re-derived by me, \textbf{[Q]} a standard fact quoted from memory. Nothing here has been
read by a human.

\section{The dictionary between sums and operators}

Let $P_n(x)=\sum_{k=0}^n(-1)^k\binom nk\binom{n+k}kx^k$ be the Legendre polynomial of the interval
$[0,1]$, normalised by $P_n(0)=1$; write $\ip fg=\int_0^1fg\,dx$, so that
$\ip{P_n}{P_m}=\frac{\delta_{nm}}{2n+1}$. The rational function of [MP] is a moment sequence,
\[
r_n(t)=\int_0^1x^{t-1}P_n(x)\,dx ,
\]
and summing over $t$ under the integral sign gives
\[
I_{n,m}=-\sum_{t\ge1}(r_nr_m)'(t)=\int_0^1\!\!\int_0^1P_n(x)\,P_m(y)\,K(xy)\,dx\,dy,\qquad
K(u)=\frac{-\log u}{1-u}.
\]
The two elementary moves of [MP] become two operators: the shift $t\to t+1$ is multiplication by
$x$, which we call $X$; and multiplication by $t$ is, up to sign, the operator $x\frac d{dx}$. So
polynomial identities in $t$ between shifted functions are identities between differential
operators. That is the whole reason operators appear.

\section{Two operators, and why the cubes}

\q{The first operator} is the one that has the $P_n$ as eigenfunctions:
\[
Jf=-\bigl(x(1-x)f'\bigr)',\qquad JP_n=\lambda_nP_n,\quad\lambda_n=n(n+1).\qquad\textbf{[V]}
\]
(In [MP] this is relation (B).) Multiplication by $x$ is tridiagonal in this basis, with weights
$n+1$ and $n$: \quad$2(2n+1)\,XP_n=-(n+1)P_{n+1}+(2n+1)P_n-nP_{n-1}$.\quad\textbf{[V]}

\q{The second operator} is the symmetrised product of the two, with a constant added:
\[
H=\tfrac12(XJ+JX)+\tfrac12,\qquad Hf=-\bigl(x^2(1-x)f'\bigr)'+xf.\qquad\textbf{[V]}
\]
Because $J$ is diagonal and $X$ is tridiagonal, $H$ is tridiagonal, and its entry between $P_n$ and
$P_{n+1}$ is the entry of $X$ times the mean of the two eigenvalues of $J$:
$(n+1)\cdot\frac12(\lambda_n+\lambda_{n+1})=(n+1)\cdot(n+1)^2$. \emph{That is where the cubes come
from in this language.} Exactly:
\begin{equation}\label{tri}
2(2n+1)\,HP_n=-(n+1)^3P_{n+1}+\bigl((n+1)^3+n^3\bigr)P_n-n^3P_{n-1}.\qquad\textbf{[V]}
\end{equation}
The constant $\frac12$ in $H$ is what makes the middle coefficient equal to the sum of the outer
ones.

\q{$H$ is the resistor chain.} Expand a function in Legendre coordinates, $\varphi_n=\ip{P_n}f$.
Then \eqref{tri} gives \textbf{[V]}
\[
\ip f{Hf}=\frac12\sum_{n\ge1}n^3\,(\varphi_n-\varphi_{n-1})^2,\qquad
\ip ff=\sum_{n\ge0}(2n+1)\,\varphi_n^2 .
\]
The first formula is the power dissipated in a chain in which node $n-1$ is joined to node $n$ by a
conductance $n^3$, with voltages $\varphi_n$. There is no leak. Adding a multiple $\lambda$ of the
identity to $H$ adds, by the second formula, a leak of conductance $2(2n+1)\lambda$ at node $n$.
Compare [MP], Section~7: the chain has resistance $\sum n^{-3}=\zeta(3)$, and indeed
\begin{equation}\label{res}
\ip1{(H+m(m+1))^{-1}1}=2\sum_{k>m}\frac1{k^3}\qquad(m=0,1,2,\dots).
\end{equation}
For $m=0$ this reads $\zeta(3)=\frac12\ip1{H^{-1}1}$. The right side of \eqref{res} is the entry
$I_{0,m}$ of the table, which I verified exactly; that the left side equals $I_{0,m}$ rests on the
identification of Section~3. An agent reports that it holds with $m$ replaced by any real
$x>-\frac12$, the right side being $2\zeta(3,x+1)$; this is Ramanujan's continued fraction, which
is the expansion of the left side along the chain. \textbf{[A]}

\section{The kernel intertwines the two operators}

The kernel of Section~1 satisfies
\begin{equation}\label{int}
\bigl(H_x+J_y\bigr)\,K(xy)=0 ,\qquad\textbf{[V]}
\end{equation}
where the subscript says on which variable the operator acts. So if $k_m(x)=\int_0^1K(xy)P_m(y)\,dy$,
then moving $J$ onto $P_m$ gives $(H+\lambda_m)\,k_m=$ a boundary term, and the boundary term is the
constant $1$. Hence
\[
I_{n,m}=\ip{P_n}{k_m}=\ip{P_n}{(H+\lambda_m)^{-1}1} .
\]
\emph{Row $m$ of the table of [MP] is the voltage along the chain with uniform leak $\lambda_m$, fed
at node $0$.} Taking the Legendre coefficients of $(H+\lambda_m)k_m=1$ with \eqref{tri} gives the
row equation (4) of [MP], together with its equation at $n=0$,
$-I_{1,m}+(1+2\lambda_m)I_{0,m}=2$.

\emph{Status.} I verified \eqref{int} symbolically, and the resulting equations for $I_{n,m}$,
including the one at $n=0$, in exact arithmetic for $n,m\le14$. The pointwise statement
$(H+\lambda_m)k_m=1$, the vanishing of the other boundary terms, and the questions of domain
($k_m$ has logarithmic singularities) were checked by an agent with series expansions and are not
proved here. \textbf{[A]} The derivation in [MP] does not depend on them.

\section{What kind of operator $H$ is}

Substitute $y=\frac2x-1$, which maps $(0,1)$ onto $(1,\infty)$, and write $f(x)=\frac1x\,g(y)$. Then
\[
Hf=\frac1x\,Lg,\qquad Lg=-\frac d{dy}\Bigl((y^2-1)\frac{dg}{dy}\Bigr),\qquad
\int_0^1f^2\,dx=\frac12\int_1^\infty g^2\,dy.\qquad\textbf{[V]}
\]
With $y=\cosh\rho$ this $L$ is the radial part of the Laplacian of the hyperbolic plane. Its
spectrum is the half-line $[\frac14,\infty)$. \textbf{[Q]} Meanwhile $J$ is the same expression on
the interval $(-1,1)$ (put $y=1-2x$): the Laplacian of the sphere, restricted to functions of the
latitude, with eigenvalues $n(n+1)$. So:
\begin{itemize}
\item $J$ lives on the sphere, $H$ on the hyperbolic plane, and \eqref{int} says that $K$ carries
one to the other with a change of sign.
\item The equation $(L+m(m+1))g=0$ is Legendre's equation of degree $m$ on $(1,\infty)$. Its two
solutions are the Legendre functions $P_m$ and $Q_m$. The numbers $I_{n,m}$ come from $Q_m$, and
their integer coefficients $\sum_j\binom nj\binom{n+j}j\binom mj\binom{m+j}j$ from $P_m$.
\textbf{[A]}
\item $H\ge\frac14$, and $\lambda_n+\frac14=(n+\frac12)^2$. This is why $n+\frac12$ is the natural
variable throughout.
\end{itemize}

\section{Sign and size in this language}

\q{The energy argument (Codex), in three lines.} There is a polynomial $w_n$ of degree $n-1$ with
$P_n=(H+\lambda_n)w_n+a_n$, where $a_n$ is Ap\'ery's integer. \textbf{[V]} for $n\le6$. Since $w_n$
is orthogonal to $P_n$,
\[
a_n\,I_n=\ip{P_n}{(H+\lambda_n)^{-1}P_n},\qquad\text{so}\qquad
0<a_n\,I_n\le\frac{\|P_n\|^2}{\lambda_n+\frac14}=\frac4{(2n+1)^3},
\]
using only that $H$ is self-adjoint and $H\ge\frac14$. (Numerically $a_nI_n(2n+1)^3$ increases to
$2\sqrt2=2.828$, so the bound is off by the factor $\sqrt2$. \textbf{[V]} numerically, $n\le59$.)
This gives sign and size at once. Its price: one needs to know how fast $a_n$ grows, and the
self-adjointness of $H$ on the right domain. \textbf{[A]}

\q{The maximum principle, in this language.} Put $\Psi=\sum_n(2n+1)I_nP_n$, the function whose
Legendre coordinates are the $I_n$. Ap\'ery's recurrence, with its equation at $n=0$, is
equivalent (coordinate by coordinate) to
\begin{equation}\label{diag}
\bigl(H+8J+2\bigr)\Psi=6,\qquad H+8J+2=-\frac d{dx}\Bigl[x(1-x)(x+8)\frac d{dx}\Bigr]+x+2 .
\qquad\textbf{[V]}
\end{equation}
In the chain picture, $8J+2$ is a leak of conductance $2(2n+1)(8\lambda_n+2)=4(2n+1)^3$ at node
$n$: the leak of [MP]. The maximum principle of [MP] says that the solution fed at node $0$ is
positive at every node and decreases along the chain, the standard behaviour of a chain with
positive conductances and positive leaks. Two things can be read off \eqref{diag}.
\begin{itemize}
\item \emph{The rate.} The equation has a singular point at $x=-8$. The Legendre coordinates of a
function decay at the rate fixed by its nearest singularity; for the interval $[0,1]$ and the
point $-8$ that rate is $17+\sqrt{17^2-1}=17+12\sqrt2$. \textbf{[Q]} for the general principle.
\item \emph{The constants.} $8J+2=6J+2(J+1)$: six times the row leak $\lambda_n$, plus four times
the diagonal entry $\frac12(\lambda_n+1)$ of $H$. The $6$ on the right of \eqref{diag} is the same
$6$.
\end{itemize}
The row picture and the diagonal picture are thus two resolvents: $(H+\lambda_m)^{-1}1$ for row
$m$, and $6\,(H+8J+2)^{-1}1$ for the diagonal.

\section{The commutators}

Let $Z=[J,X]=JX-XJ$. Then \textbf{[V]}
\[
Z=-2x(1-x)\frac d{dx}-(1-2x),\qquad [J,Z]=2(JX+XJ)-2J,\qquad [Z,X]=2X^2-2X .
\]
So $J$, $X$ and their commutator close under commutation up to quadratic terms. This is the
algebraic form of the statement that each of $J$ and $X$ is tridiagonal in the eigenbasis of the
other. The pair $(J,H)$ satisfies relations of the same kind, and is symmetric under
$J\leftrightarrow-H$, which is the symmetry $n\leftrightarrow m$ of the table. \textbf{[A]} From
memory, such pairs are called Leonard pairs, and the passage from $(J,X)$ to $H$ is the
``tridiagonalisation'' of Ismail and Koelink. \textbf{[Q]}

\q{What the algebra buys.} Two relations between the four entries around a square of the table,
valid for all $n,m\ge1$, for the numbers $I_{n,m}$ and for their integer coefficients:
\begin{align*}
(n-m)\,\bigl(X_{n,m}-X_{n-1,m-1}\bigr)&=(n+m)\,\bigl(X_{n-1,m}-X_{n,m-1}\bigr),\\
(n^2-nm+m^2)\,\bigl(X_{n,m}+X_{n-1,m-1}\bigr)&=(n^2+nm+m^2)\,\bigl(X_{n,m-1}+X_{n-1,m}\bigr).
\end{align*}
\textbf{[V]} exactly for $n,m\le10$. For $m=n$ the second one is $I_n+I_{n-1}=6\,I_{n,n-1}$: the $6$
is $2\cdot\frac{3n^2}{n^2}$. Applied at $(n+1,n)$ and combined, the two relations give Ap\'ery's
recurrence without the row equation. \textbf{[A]} So the table carries more structure than [MP]
uses.

\q{The same scheme for the other constants.} The three classical cases use three operators built
from $J$ and $X$ of degree $0$, $1$, $2$ in $J$:
\begin{center}
\begin{tabular}{@{}llll@{}}
\toprule
constant & operator & in the Legendre basis & weights\\
\midrule
$\log2$ & $X$ & tridiagonal, symmetric & $n+1,\ n$\\
$\zeta(2)$ & $\frac12(1+[J,X])=x-x(1-x)\frac d{dx}$ & tridiagonal, not symmetric & $-(n+1)^2,\ +n^2$\\
$\zeta(3)$ & $\frac12(XJ+JX+1)=H$ & tridiagonal, symmetric & $(n+1)^3,\ n^3$\\
\bottomrule
\end{tabular}
\end{center}
The tridiagonal forms are \textbf{[V]}; that the corresponding tables give the classical
recurrences for $\log2$ and $\zeta(2)$ is \textbf{[A]} (exact for $n,m\le12$). The commutator is
antisymmetric, which is why the weights for $\zeta(2)$ have opposite signs and its small solution
alternates. The agents did not find an operator of this kind for $\zeta(5)$ with positive weights.
\textbf{[A]}; a negative search result, not a theorem.

\section{What the language buys, and what it does not}

\emph{It buys understanding.} The cubes are (entry of $X$) times (mean eigenvalue of $J$). The
leak-free chain is $H$ itself, and $\zeta(3)$ is its resistance. Rows and diagonal are resolvents.
The number $\frac14$ and the variable $n+\frac12$ come from hyperbolic geometry. The rate
$17+12\sqrt2$ is the position of a singular point. The three constants are one construction.

\emph{It buys one short argument:} the energy bound of Section~5.

\emph{It does not buy the arithmetic.} That $a_n$ is an integer and that $D_n^3b_n$ is an integer
still come from the explicit formulas. Nor does it prove the size without either the growth of
$a_n$ or the recurrence.

\emph{It costs rigour.} Statements about $H$ as a self-adjoint operator need its domain: an agent
reports that $H$ is in the limit-point case at $x=0$ and the limit-circle case at $x=1$ (with
the boundary condition $(1-x)f'\to0$), and that the standard criterion for the uniqueness of the
self-adjoint extension from the tridiagonal form fails, so a separate argument is needed.
\textbf{[A]} The proof in [MP] avoids all of this by working with sums.

\emph{Relation to the literature} (from the agents' search, not re-read by me). The operator $H$
is the tridiagonalisation operator of Ismail and Koelink (arXiv:1108.5716) and an ``algebraic Heun
operator'' in the sense of Genest, Ismail, Vinet and Zhedanov (arXiv:1506.07803). The agents found
no paper that presents Ap\'ery's proof through a self-adjoint operator, a Green's function or an
energy. That is a statement about a short search, not a claim of novelty.
\end{document}
